% ---------------------------------------------------------------------------
% Shared preamble: packages, theorem environments, notation macros, and the
% listings style for delhi source.
% ---------------------------------------------------------------------------
\usepackage[T1]{fontenc}
\usepackage[utf8]{inputenc}
\usepackage[english]{babel}
\usepackage{lmodern}
\usepackage{microtype}

\usepackage{amsmath, amssymb, amsthm}
\usepackage{mathtools}
\usepackage{booktabs}
\usepackage{array}
\usepackage{listings}
\usepackage{xcolor}
\usepackage{tikz}
\usetikzlibrary{arrows.meta, positioning, fit, backgrounds, calc, shapes.geometric}
\usepackage[round,authoryear]{natbib}

% --- model-diagram styles --------------------------------------------------
% Worlds are circles labelled with the atoms true in them; the designated world
% carries a double ring. Plausibility edges point toward the *more* plausible world.
\tikzset{
  world/.style        = {circle, draw, minimum size=9.5mm, inner sep=1pt, font=\small},
  actual/.style       = {world, double, double distance=0.6pt},
  event/.style        = {rectangle, rounded corners=2pt, draw, minimum height=7mm,
                         minimum width=11mm, inner sep=2pt, font=\small},
  plaus/.style        = {-{Stealth[length=2mm]}, shorten >=1pt, shorten <=1pt},
  equi/.style         = {{Stealth[length=2mm]}-{Stealth[length=2mm]}, shorten >=1pt,
                         shorten <=1pt},
  elab/.style         = {font=\scriptsize, inner sep=1.5pt, fill=white},
  wlab/.style         = {font=\scriptsize\itshape, inner sep=1pt},
}
\usepackage[hidelinks]{hyperref}
\usepackage{cleveref}

% --- theorem environments --------------------------------------------------
% All theorem-like environments share one numbering sequence, so a reader meets
% Definition 2.1, Proposition 2.2, Theorem 2.3 rather than three interleaved series.
%
% The obvious way to do that -- \newtheorem{proposition}[definition]{Proposition} -- makes
% every environment literally use the counter named `definition`, and cleveref names a
% reference after its counter. The result is that \Cref of a proposition renders
% "Definition 2.5". `aliascnt` gives each environment a distinct counter name that is an
% alias for the same underlying counter, which keeps the shared sequence and lets cleveref
% tell them apart.
% All theorem-like environments share the `definition` counter, so a reader meets
% Definition 2.1, Proposition 2.2, Theorem 2.3 in one sequence.
%
% A consequence: cleveref names a reference after its *counter*, so \Cref of a
% proposition would render "Definition 2.5". The `aliascnt` package is the usual remedy,
% but it sends this document into an unterminating loop during AMS symbol-font setup.
% Cross-references to non-definition environments are therefore written explicitly --
% "Proposition~\ref{...}" rather than "\Cref{...}" -- and \cref is reserved for
% definitions, sections, tables and equations, where it names them correctly.
\theoremstyle{definition}
\newtheorem{definition}{Definition}[section]
\newtheorem{example}[definition]{Example}
\newtheorem{remark}[definition]{Remark}

\theoremstyle{plain}
\newtheorem{proposition}[definition]{Proposition}
\newtheorem{theorem}[definition]{Theorem}
\newtheorem{lemma}[definition]{Lemma}
\newtheorem{corollary}[definition]{Corollary}

% --- notation --------------------------------------------------------------
% Modalities are written with an explicit agent subscript throughout.
\newcommand{\Ag}{\mathcal{G}}                  % set of agents
\newcommand{\Props}{\mathcal{P}}               % set of propositions
\newcommand{\Worlds}{W}
\newcommand{\Kop}[1]{K_{#1}}                   % knowledge
\newcommand{\Bop}[1]{B_{#1}}                   % belief
\newcommand{\Sop}[1]{\Box_{#1}}                % safe belief
\newcommand{\Cop}[1]{C_{#1}}                   % common knowledge
\newcommand{\CBop}[2]{B^{#2}_{#1}}             % conditional belief
\newcommand{\Kw}[1]{\mathit{Kw}_{#1}}
\newcommand{\Bw}[1]{\mathit{Bw}_{#1}}
\newcommand{\Ig}[1]{\mathord{?}_{#1}}          % ignorance
\newcommand{\Susp}[1]{\mathord{\text{\textquestiondown}}_{#1}}

\newcommand{\plaus}[1]{R_{#1}}                 % plausibility preorder
\newcommand{\comp}[1]{\sim_{#1}}               % comparability
\newcommand{\best}[1]{\rightarrow_{#1}}        % maxima operator
\newcommand{\Lang}{\mathcal{L}}
\newcommand{\Lgb}{\Lang^{\Props}_{\Ag}}        % full modal language
\newcommand{\Lprop}{\Lang^{\Props}}            % propositional fragment
% Semantic brackets, defined here rather than via stmaryrd so the paper builds on a
% minimal TeX installation.
\providecommand{\llbracket}{[\![}
\providecommand{\rrbracket}{]\!]}
\newcommand{\denote}[1]{\llbracket #1 \rrbracket}
\newcommand{\bisimR}{\mathbin{\underline{\leftrightarrow}_{R}}}
\newcommand{\bisimD}{\mathbin{\underline{\leftrightarrow}_{D}}}
\newcommand{\ML}{\mathit{ML}}
\newcommand{\pre}{\mathit{pre}}
\newcommand{\add}{\mathit{add}}
\newcommand{\del}{\mathit{del}}
\newcommand{\FPN}{\mathit{FPN}}
\newcommand{\PN}{\mathit{PN}}
\newcommand{\Nlab}{\mathit{N}}
\newcommand{\hole}{\underline{\phantom{x}}}

% --- listings style for delhi source ---------------------------------------
\definecolor{kwcolor}{HTML}{6A3D9A}
\definecolor{clcolor}{HTML}{B15928}
\definecolor{cmcolor}{HTML}{6B7280}
\definecolor{opcolor}{HTML}{1F78B4}
\definecolor{bgcolor}{HTML}{F7F7F9}

\lstdefinelanguage{delhi}{
  morekeywords={types,objects,agents,props,constants,define,rules,initially,
                state,goal,invariants,actions},
  morekeywords=[2]{actor,pre,causes,determines,announces,observes,aware,if},
  morekeywords=[3]{K,B,C,Kw,Bw},
  sensitive=true,
  morecomment=[l]{//},
  morecomment=[s]{/*}{*/},
  morestring=[b]",
}

\lstset{
  language=delhi,
  basicstyle=\ttfamily\small,
  keywordstyle=\color{kwcolor}\bfseries,
  keywordstyle=[2]\color{clcolor},
  keywordstyle=[3]\color{opcolor},
  commentstyle=\color{cmcolor}\itshape,
  backgroundcolor=\color{bgcolor},
  frame=single,
  rulecolor=\color{gray!40},
  framesep=5pt,
  xleftmargin=10pt,
  xrightmargin=4pt,
  breaklines=true,
  showstringspaces=false,
  columns=fullflexible,
  keepspaces=true,
  captionpos=b,
  aboveskip=1.1em,
  belowskip=0.9em,
}

% Short listings are discussed in the paragraph immediately after them, so splitting one
% across a page break strands a lone brace with the caption. `keepit` wraps a listing in a
% minipage, which cannot break. Use it for anything under about fifteen lines; longer
% listings should be floats instead.
\newenvironment{keepit}
  {\par\medskip\noindent\begin{minipage}{\linewidth}}
  {\end{minipage}\par\medskip}

% A narrower style for shell transcripts.
\lstdefinestyle{shell}{
  language={},
  basicstyle=\ttfamily\footnotesize,
  keywordstyle={},
  commentstyle=\color{cmcolor}\itshape,
  morecomment=[l]{\#},
}

% --- misc ------------------------------------------------------------------
\newcommand{\delhi}{\textsf{delhi}}
\newcommand{\mB}{\textsf{mB}}
\newcommand{\mBplus}{\textsf{mB${}^{+}$}}
\newcommand{\code}[1]{\texttt{\small #1}}
